\documentclass[aspectratio=169, 10pt]{beamer}
\usetheme{metropolis}

% --- Edinburgh colours ----------------------------------------
\definecolor{UoERed}{HTML}{7A2318}
\definecolor{UoEGold}{HTML}{B8860B}
\definecolor{CodeBG}{HTML}{F7F3ED}
\definecolor{CodeFrame}{HTML}{D5C9B1}

\setbeamercolor{palette primary}{bg=UoERed, fg=white}
\setbeamercolor{frametitle}{bg=UoERed, fg=white}
\setbeamercolor{title separator}{fg=UoEGold}
\setbeamercolor{progress bar}{fg=UoEGold, bg=UoEGold!20}
\setbeamercolor{alerted text}{fg=UoERed}
\setbeamercolor{block title}{bg=UoERed!15, fg=UoERed}
\setbeamercolor{block body}{bg=UoERed!5}

% --- Packages -------------------------------------------------
\usepackage[T1]{fontenc}
\usepackage{lmodern}
\usepackage{amsmath, amssymb, mathtools}
\usepackage{listings}
\usepackage{graphicx, booktabs}
\usepackage{tikz}
\usetikzlibrary{arrows.meta, positioning, calc, decorations.pathreplacing}
\usepackage{hyperref}
\hypersetup{colorlinks=true, linkcolor=UoERed, urlcolor=UoERed!70!black}
\setbeamertemplate{navigation symbols}{}

\lstset{
  language=Python,
  basicstyle=\ttfamily\scriptsize,
  keywordstyle=\color{UoERed}\bfseries,
  commentstyle=\color{gray}\itshape,
  stringstyle=\color{UoEGold},
  backgroundcolor=\color{CodeBG},
  frame=single,
  rulecolor=\color{CodeFrame},
  numbers=none, breaklines=true, showstringspaces=false, tabsize=4,
  xleftmargin=4pt, xrightmargin=4pt, aboveskip=6pt, belowskip=4pt,
}

\AtBeginSection[]{
  \begin{frame}
    \vfill\centering
    \usebeamerfont{title}\insertsectionhead\par
    \vfill
  \end{frame}
}

\title{Week 5 --- Constrained Optimisation II}
\subtitle{KKT Conditions \& Inequality Constraints}
\author{Dr Juan Zurita}
\institute{Optimisation \& Mathematical Methods for Economics\\
The University of Edinburgh~$\cdot$~School of Economics}
\date{}

\begin{document}
\maketitle

\begin{frame}{Inequality Constraints}
Most economic constraints are \textbf{inequalities}: budgets, non-negativity, capacity.\\\bigskip $\min f(\mathbf{x})$ s.t.\ $g_i(\mathbf{x}) \leq 0$, $i=1,\ldots,m$
\end{frame}

\begin{frame}{KKT Conditions}
\begin{enumerate}\item \textbf{Stationarity:} $\nabla f + \sum\mu_i\nabla g_i = 0$\item \textbf{Primal feasibility:} $g_i(\mathbf{x}) \leq 0$\item \textbf{Dual feasibility:} $\mu_i \geq 0$\item \textbf{Complementary slackness:} $\mu_i\,g_i(\mathbf{x}) = 0$\end{enumerate}\par\bigskip Under convexity, KKT conditions are \alert{sufficient}.
\end{frame}

\begin{frame}{Complementary Slackness}
$\mu_i\,g_i = 0$ means:\\\bigskip\begin{itemize}\item Constraint \textbf{binding} ($g_i = 0$): multiplier can be positive\item Constraint \textbf{slack} ($g_i < 0$): multiplier must be zero\end{itemize}\par\bigskip\textbf{Economics:} a resource has a positive shadow price only if it is fully used.
\end{frame}

\begin{frame}{Portfolio Choice}
$\min\;\mathbf{w}^T\Sigma\mathbf{w}$ s.t.\ $\mathbf{w}^T\boldsymbol{\mu} \geq r_{\min}$, $\sum w_i = 1$, $w_i \geq 0$\\\bigskip No short-selling constraints are \textbf{inequalities}.\\\bigskip KKT tells us: assets with zero weight have insufficient expected return to justify inclusion.
\end{frame}

\begin{frame}{Summary}
\begin{itemize}\item KKT = Lagrange multipliers extended to inequalities\item Complementary slackness: binding vs.\ slack constraints\item Under convexity, KKT is necessary \textit{and} sufficient\item \textbf{Next week:} convex optimisation theory and CVXPY\end{itemize}
\end{frame}
\end{document}
